% !TEX root = ../main.tex
%% Chapitre 22 — Listes
%% Sources : frontend/cours/resume/terminal/resume_algorithme_listes_terminal.txt
%%           frontend/exercices/terminal_spé_Maths/exercice_algorithme_listes.txt
%% Tous les programmes ont été exécutés : les affichages indiqués sont exacts.
%% Exercices 6 à 10 : exercices inédits rédigés pour le livre.

\begin{ouverture}[listes]{22}{Listes}{Partie 4 · Algorithmique et logique}
  {Création\enspace·\enspace Indices\enspace·\enspace Méthodes\enspace·\enspace Parcours\enspace·\enspace
   Listes en compréhension}
  \objectif{Créer une liste et accéder à ses éléments par leur indice.}
  \objectif{Modifier une liste à l'aide de ses méthodes : ajout, insertion, suppression, tri.}
  \objectif{Parcourir une liste par ses valeurs ou par ses indices.}
  \objectif{Construire une liste par ajouts successifs ou en compréhension.}
  \objectif{Utiliser des listes pour traiter une série statistique, trier ou simuler.}
  \prerequis{Boucles \code{for} et \code{while}, accumulateurs (chapitre 21).}
  \prerequis{Statistiques descriptives : moyenne, médiane, écart-type.}
\end{ouverture}
\miseEnPageSouple

%% =====================================================================
%%  L'ESSENTIEL
%% =====================================================================
\essentiel

\essSection{1}{Créer une liste et accéder à ses éléments}

\begin{definition}[Liste]
Une \textbf{liste} est une suite \emph{ordonnée} et \emph{modifiable} de valeurs, écrites entre crochets et séparées par des
virgules, par exemple \code{L = [12, 15, 18, 10, 14]}. Chaque élément est repéré par son \textbf{indice} : le premier a l'indice $0$, le dernier l'indice \code{len(L) - 1}, aussi noté $-1$.
\end{definition}

\begin{center}
\begin{tikzpicture}[x=12mm,y=1mm]
  \foreach \v [count=\i from 0] in {12,15,18,10,14} {
    \draw[otNavy,fill=otNavyTint,line width=0.6pt] (\i,0) rectangle ++(1,8);
    \node at (\i+0.5,4) {$\v$};
    \node[font=\otSemi\fontsize{8}{9}\selectfont,text=otNavy] at (\i+0.5,11) {\i};
    \pgfmathtruncatemacro\neg{\i-5}
    \node[font=\fontsize{8}{9}\selectfont,text=otMuted] at (\i+0.5,-3.2) {$\neg$};
  }
  \node[anchor=east,font=\small] at (-0.15,4) {\code{L =}};
  \node[anchor=west,font=\scriptsize,text=otNavy] at (5.2,11) {indices positifs};
  \node[anchor=west,font=\scriptsize,text=otMuted] at (5.2,-3.2) {indices négatifs};
  \draw[otGreenDark,line width=1.1pt,rounded corners=1mm] (1.06,0.9) rectangle (2.94,7.1);
  \node[anchor=north,font=\scriptsize,text=otGreenDark] at (2,-6.5) {\code{L[1:3]} $\to$ \code{[15, 18]}};
\end{tikzpicture}
\end{center}

\noindent
\begin{minipage}[t]{0.48\linewidth}\vspace{0pt}
{\renewcommand{\arraystretch}{1.3}\arrayrulecolor{otRule}%
\begin{tabularx}{\linewidth}{@{}lYl@{}}
\toprule
\otTete{Instruction} & \otTete{Rôle} & \otTete{Résultat} \\
\midrule
\code{L[0]} & premier élément & \code{12} \\
\code{L[-1]} & dernier élément & \code{14} \\
\code{len(L)} & nombre d'éléments & \code{5} \\
\code{L[1:3]} & indices $1$ et $2$ & \code{[15, 18]} \\
\bottomrule
\end{tabularx}}
\end{minipage}\hfill
\begin{minipage}[t]{0.47\linewidth}\vspace{0pt}
\begin{propriete}[Modifier un élément]
\raggedright
\code{L[1] = 16} remplace l'élément d'indice $1$ : \code{L} devient \code{[12, 16, 18, 10, 14]}. Un indice qui n'existe pas, comme
\code{L[5]}, provoque une erreur \code{IndexError}.
\end{propriete}
\end{minipage}

\essSection{2}{Méthodes et fonctions utiles}

\noindent
\begin{minipage}[t]{0.55\linewidth}\vspace{0pt}
{\renewcommand{\arraystretch}{1.3}\arrayrulecolor{otRule}%
\begin{tabularx}{\linewidth}{@{}lYl@{}}
\toprule
\otTete{Méthode} & \otTete{Effet} & \otTete{\code{L} devient} \\
\midrule
\code{L.append(5)} & ajoute à la fin & \code{[3, 1, 2, 5]} \\
\code{L.insert(0, 7)} & insère à l'indice $0$ & \code{[7, 3, 1, 2, 5]} \\
\code{L.remove(1)} & supprime la 1\textsuperscript{re} occurrence & \code{[7, 3, 2, 5]} \\
\code{L.pop()} & retire et renvoie le dernier & \code{[7, 3, 2]} \\
\code{L.sort()} & trie dans l'ordre croissant & \code{[2, 3, 7]} \\
\code{L.reverse()} & inverse l'ordre & \code{[7, 3, 2]} \\
\bottomrule
\end{tabularx}}\\[1mm]
\legende{Méthodes appliquées successivement à \code{L = [3, 1, 2]}.}
\end{minipage}\hfill
\begin{minipage}[t]{0.41\linewidth}\vspace{0pt}
{\renewcommand{\arraystretch}{1.3}\arrayrulecolor{otRule}%
\begin{tabularx}{\linewidth}{@{}lY@{}}
\toprule
\otTete{Fonction} & \otTete{Renvoie} \\
\midrule
\code{sum(L)} & la somme des éléments \\
\code{min(L)} & le plus petit élément \\
\code{max(L)} & le plus grand élément \\
\code{sorted(L)} & une \emph{copie} triée \\
\code{x in L} & \code{True} si \code{x} est dans \code{L} \\
\code{L.index(x)} & l'indice de la 1\textsuperscript{re} occurrence \\
\code{L.count(x)} & le nombre d'occurrences \\
\bottomrule
\end{tabularx}}
\end{minipage}

\essSection{3}{Parcourir une liste}

\begin{codecote}[0.52]
\begin{pythoncol}
fruits = ["pomme", "banane", "orange"]
for f in fruits:
    print(f, end=" ")
\end{pythoncol}

\smallskip
\begin{consolecol}
pomme banane orange
\end{consolecol}
\cote
\begin{pythoncol}
for i in range(len(fruits)):
    print(i, fruits[i])
\end{pythoncol}

\smallskip
\begin{consolecol}
0 pomme
1 banane
2 orange
\end{consolecol}
\end{codecote}

\begin{propriete}[Deux façons de parcourir]
Le parcours \textbf{par valeurs} (\code{for x in L:}) suffit quand seules les valeurs comptent ; le parcours \textbf{par indices}
(\code{for i in range(len(L)):}) est nécessaire quand on a besoin de la position, par exemple pour modifier \code{L[i]} ou comparer
deux éléments voisins.
\end{propriete}

\essSection{4}{Construire une liste}

\begin{savoirfaire}[Construire la liste des termes d'une suite]
\begin{codecote}[0.44]
\begin{pythoncol}
u = [3]
for k in range(5):
    u.append(2 * u[-1] - 1)
print(u)
\end{pythoncol}
\cote
On part de la liste \code{[3]} contenant $u_0$ ; à chaque tour, on calcule le terme suivant à partir du dernier, \code{u[-1]}, et on
l'ajoute avec \code{append}.
\begin{consolecol}
[3, 5, 9, 17, 33, 65]
\end{consolecol}
\end{codecote}
\end{savoirfaire}

\begin{definition}[Liste en compréhension]
L'écriture \code{[expression for x in …]} construit en une ligne la liste des valeurs de l'expression ; on peut ajouter un filtre
\code{if condition}.\\[1mm]
\code{[k**2 for k in range(1, 6)]} $\to$ \code{[1, 4, 9, 16, 25]}\\
\code{[x for x in [15, 8, 18, 10, 16] if x >= 12]} $\to$ \code{[15, 18, 16]}
\end{definition}

\begin{exemple}[Une affectation ne copie pas une liste]
\begin{codecote}[0.4]
\begin{pythoncol}
L = [1, 2, 3]
M = L
M.append(4)
print(L)
N = L.copy()
N.append(5)
print(L, N)
\end{pythoncol}
\cote
\begin{consolecol}
[1, 2, 3, 4]
[1, 2, 3, 4] [1, 2, 3, 4, 5]
\end{consolecol}

\smallskip
\centering
\begin{tikzpicture}[x=7mm,y=7mm]
  \foreach \v [count=\i from 0] in {1,2,3,4} {
    \draw[otNavy,fill=otNavyTint] (\i,0) rectangle ++(1,1); \node at (\i+0.5,0.5) {\small$\v$}; }
  \foreach \v [count=\i from 0] in {1,2,3,4,5} {
    \draw[otGreenDark,fill=otGreenTint] (\i,-2) rectangle ++(1,1); \node at (\i+0.5,-1.5) {\small$\v$}; }
  \tikzset{nomvar/.style={draw=otMuted,fill=white,rounded corners=0.8mm,minimum width=6.5mm,minimum height=5mm,inner sep=0pt}}
  \node[nomvar] (l) at (-2.2,1.05) {\code{L}};
  \node[nomvar] (m) at (-2.2,-0.05) {\code{M}};
  \node[nomvar] (n) at (-2.2,-1.5) {\code{N}};
  \draw[otFleche] (l.east) -- (-0.08,0.65);
  \draw[otFleche] (m.east) -- (-0.08,0.35);
  \draw[otFleche] (n.east) -- (-0.08,-1.5);
\end{tikzpicture}\\[0.5mm]
\legende{\code{L} et \code{M} désignent la même liste ; \code{N} est une copie.}
\end{codecote}
\end{exemple}

\begin{attention}[Pièges fréquents]
\begin{itemize}[label={\color{otOrange}\textbullet},leftmargin=1.4em,itemsep=0.6mm,topsep=0.5mm]
  \item Le premier élément a l'indice $0$ : une liste de $n$ éléments a pour indices $0$, $1$, …, $n-1$.
  \item \code{L.sort()} modifie \code{L} et ne renvoie rien : écrire \code{L = L.sort()} remplace la liste par \code{None}. Pour obtenir
    une copie triée, on utilise \code{sorted(L)}.
  \item \code{M = L} ne copie pas la liste : \code{M} et \code{L} désignent la \emph{même} liste, et modifier l'une modifie l'autre.
    Pour une copie indépendante : \code{M = L.copy()}.
  \item \code{L.remove(x)} et \code{L.index(x)} provoquent une erreur si \code{x} n'est pas dans la liste : on teste d'abord
    \code{if x in L:}.
\end{itemize}
\end{attention}

%% =====================================================================
%%  EXERCICES
%% =====================================================================
\exercices
\rubrique{Découverte}

\begin{exercice}[1]{Créer une liste et accéder à ses éléments}
\begin{questions}
  \item On pose \code{notes = [12, 15, 18, 10, 14]}. Écrire les instructions qui affichent la première note, la dernière note et la
    note du milieu.
  \item On pose \code{fruits = ["pomme", "banane", "orange", "kiwi"]}. Donner la valeur de :
\begin{tasks}[label=\alph*)](4)
  \task \code{fruits[2]}
  \task \code{fruits[-2]}
  \task \code{fruits[1:3]}
  \task \code{len(fruits)}
\end{tasks}
  \item Que se passe-t-il si l'on demande \code{fruits[4]} ? Expliquer.
\end{questions}
\end{exercice}

\begin{exercice}[1]{Modifier une liste}
\begin{questions}
  \item Partir d'une liste vide \code{nombres = []}, y ajouter $5$, $10$ et $15$ avec \code{append}, puis insérer $8$ à l'indice $1$
    avec \code{insert}. Donner le contenu de la liste après chaque étape.
  \item Partir de \code{couleurs = ["rouge", "bleu", "vert", "bleu"]}, supprimer le premier \code{"bleu"} avec \code{remove}, puis
    remplacer \code{"vert"} par \code{"jaune"}.
  \item On exécute ensuite \code{x = couleurs.pop()}. Que valent \code{x} et \code{couleurs} ?
\end{questions}
\end{exercice}

\begin{exercice}[1]{Rechercher dans une liste}
\begin{questions}
  \item On pose \code{notes = [15, 12, 18, 12, 10]}. Donner la valeur de :
\begin{tasks}[label=\alph*)](2)
  \task \code{len(notes)}
  \task \code{notes.index(12)}
  \task \code{notes.count(12)}
  \task \code{13 in notes}
\end{tasks}
  \item On pose \code{eleves = ["Alice", "Bob", "Charlie", "Alice"]}. Donner la valeur de :
\begin{tasks}[label=\alph*)](3)
  \task \code{"Alice" in eleves}
  \task \code{eleves.index("Alice")}
  \task \code{eleves.count("Alice")}
\end{tasks}
  \item Pourquoi l'instruction \code{notes.index(13)} provoque-t-elle une erreur ? Écrire quelques lignes qui affichent la position de
    $13$ s'il est dans la liste, et « absent » sinon.
\end{questions}
\end{exercice}

\rubrique{Entraînement}

\begin{exercice}[2]{Trier une liste}
\begin{questions}
  \item Avec \code{notes = [15, 8, 20, 12]}, trier la liste avec \code{sort}, puis l'inverser avec \code{reverse}. Afficher la liste à
    chaque étape.
  \item Trier la liste \code{prenoms = ["Charlie", "Alice", "Bob"]} par ordre alphabétique, puis afficher le premier et le dernier
    prénom.
  \item Qu'affiche le programme suivant ? Quelle différence y a-t-il entre \code{sorted(L)} et \code{L.sort()} ?
\begin{python}
L = [15, 8, 20, 12]
M = sorted(L)
print(L)
print(M)
\end{python}
  \item Pourquoi l'instruction \code{L = L.sort()} est-elle une erreur ?
\end{questions}
\end{exercice}

\begin{exercice}[2]{Parcourir une liste}
\begin{questions}
  \item Avec \code{notes = [14, 12, 16, 10, 18]}, calculer la somme et la moyenne des notes à l'aide d'une boucle \code{for}.
  \item Avec \code{nombres = [5, 8, 12, 3, 15]}, déterminer le plus grand nombre à l'aide d'une boucle, sans utiliser \code{max}.
  \item Écrire un programme qui affiche le plus petit élément de \code{nombres} et son indice, sans utiliser \code{min}.
\end{questions}
\end{exercice}

\begin{exercice}[2]{Listes en compréhension}
\begin{questions}
  \item Construire en compréhension :
\begin{tasks}[label=\alph*)](2)
  \task les carrés des entiers de $1$ à $10$ ;
  \task les multiples de $7$ inférieurs à $100$ ;
  \task les termes $u_0$, …, $u_9$ de la suite $u_n=3n+2$ ;
  \task les diviseurs positifs de $36$.
\end{tasks}
  \item On pose \code{xs = [k / 4 for k in range(9)]} et \code{ys = [x**2 - 2*x for x in xs]}. Donner \code{xs} et \code{ys}, puis
    placer les points $(x\pv y)$ dans un repère. Quelle courbe obtient-on ? Quel est le minimum de la fonction sur $\intFF{0}{2}$ ?
  \item Calculer la somme $1^{2}+2^{2}+\dots+10^{2}$ avec \code{sum} et une liste en compréhension. Vérifier le résultat avec la
    formule $\frac{n(n+1)(2n+1)}{6}$.
\end{questions}
\end{exercice}

\rubrique{Problèmes}

\begin{exercice}[2]{Statistiques d'une série de notes}
\begin{codecote}[0.52]
\raggedright
Un élève a obtenu les notes suivantes ce trimestre :\\[1mm]
\code{notes = [12, 15, 8, 17, 10, 14, 15, 9]}\\[1mm]
On veut calculer les indicateurs de la série à l'aide de listes, sans fonction statistique toute faite.
\cote
\centering
\begin{tikzpicture}
\begin{axis}[otSchema, width=64mm, height=34mm, ybar, bar width=2.4mm, bar shift=0pt, xmin=7, xmax=18, ymin=0, ymax=2.6,
  xtick={8,...,17}, ytick={1,2}, xlabel={note}, xlabel style={at={(1,0)},anchor=north east,yshift=-3.5mm}]
  \addplot[fill=otNavySoft,draw=otNavy,line width=0.4pt] coordinates {(8,1) (9,1) (10,1) (12,1) (14,1) (15,2) (17,1)};
\end{axis}
\end{tikzpicture}\\[0.5mm]
\legende{Effectif de chaque note.}
\end{codecote}
\begin{questions}
  \item Calculer la moyenne $m$ à l'aide de \code{sum} et \code{len}.
  \item La variance est $V=\frac1n\sum(x_i-m)^{2}$. Calculer $V$ à l'aide d'une liste en compréhension, puis l'écart-type
    $\sigma=\sqrt V$.
  \item Écrire une fonction \code{mediane(L)} qui renvoie la médiane d'une liste de nombres, que son effectif soit pair ou impair.
    Calculer la médiane des notes.
  \item Calculer l'étendue de la série et la proportion de notes supérieures ou égales à $10$.
\end{questions}
\end{exercice}

\begin{exercice}[3]{Le tri par sélection}
\begin{codecote}[0.55]
\begin{pythoncol}
def tri_selection(L):
    n = len(L)
    for i in range(n - 1):
        imin = i
        for j in range(i + 1, n):
            if L[j] < L[imin]:
                imin = j
        L[i], L[imin] = L[imin], L[i]
    return L
\end{pythoncol}
\cote
À chaque tour de la boucle extérieure, on cherche le plus petit élément de la partie non triée, \code{L[i]}, …, \code{L[n - 1]}, et on
l'échange avec \code{L[i]}.\\[1mm]
C'est l'un des algorithmes de tri les plus simples ; il sert de référence pour comparer des méthodes plus rapides.
\end{codecote}
\begin{questions}
  \item Donner le contenu de la liste à la fin de chaque tour de la boucle extérieure lors de l'appel
    \code{tri\_selection([5, 2, 8, 1, 9])}.
  \item Expliquer pourquoi, à la fin du tour d'indice $i$, les éléments \code{L[0]}, …, \code{L[i]} sont rangés dans l'ordre
    croissant et à leur place définitive.
  \item Montrer que, pour une liste de $n$ éléments, la comparaison \code{L[j] < L[imin]} est effectuée $\frac{n(n-1)}{2}$ fois.
    Calculer ce nombre pour $n=5$ et pour $n=1\,000$.
  \item Par combien est multiplié, environ, le nombre de comparaisons lorsque la taille de la liste double ?
\end{questions}
\end{exercice}

\begin{exercice}[3]{Le crible d'Ératosthène}
Pour obtenir tous les nombres premiers jusqu'à $n$, on écrit les entiers de $2$ à $n$, puis on barre les multiples de chaque nombre
non barré, en commençant par $2$.
\begin{python}
def crible(n):
    est = [True] * (n + 1)
    est[0] = False
    est[1] = False
    for i in range(2, n + 1):
        if est[i]:
            for j in range(2 * i, n + 1, i):
                est[j] = False
    return [k for k in range(n + 1) if est[k]]
\end{python}
\begin{codecote}[0.55]
\begin{questions}
  \item Que contient la liste \code{est} au départ ? Que signifie \code{est[k] == True} à la fin du programme ?
  \item Pourquoi la boucle intérieure commence-t-elle à \code{2 * i} ? Quels nombres sont barrés lors du tour $i=3$ ?
  \item Que renvoie \code{crible(50)} ? Combien y a-t-il de nombres premiers inférieurs ou égaux à $100$ ? à $1\,000$ ?
  \item Justifier qu'il suffit d'effectuer les tours pour lesquels $i^{2}\leq n$, et que la boucle intérieure peut commencer à
    \code{i * i}.
\end{questions}
\cote
\centering
\scalebox{0.86}{\begin{tikzpicture}[x=7.2mm,y=-7.2mm]
\node[draw=otRule,fill=otNavyTint!45,text=otMuted!70,minimum size=7.2mm,inner sep=0pt,font=\fontsize{8}{9}\selectfont] at (0,0) {1};
\node[draw=otRule,fill=otGreenTint,text=otGreenDark,minimum size=7.2mm,inner sep=0pt,font=\otSemi\fontsize{8}{9}\selectfont] at (1,0) {2};
\node[draw=otRule,fill=white,text=otInk,minimum size=7.2mm,inner sep=0pt,font=\fontsize{8}{9}\selectfont] at (2,0) {3};
\node[draw=otRule,fill=otNavyTint!45,text=otMuted!70,minimum size=7.2mm,inner sep=0pt,font=\fontsize{8}{9}\selectfont] at (3,0) {4};
\draw[otMuted!80,line width=0.5pt] (2.7199999999999998,0.28) -- (3.2800000000000002,-0.28);
\node[draw=otRule,fill=white,text=otInk,minimum size=7.2mm,inner sep=0pt,font=\fontsize{8}{9}\selectfont] at (4,0) {5};
\node[draw=otRule,fill=otNavyTint!45,text=otMuted!70,minimum size=7.2mm,inner sep=0pt,font=\fontsize{8}{9}\selectfont] at (5,0) {6};
\draw[otMuted!80,line width=0.5pt] (4.72,0.28) -- (5.28,-0.28);
\node[draw=otRule,fill=white,text=otInk,minimum size=7.2mm,inner sep=0pt,font=\fontsize{8}{9}\selectfont] at (6,0) {7};
\node[draw=otRule,fill=otNavyTint!45,text=otMuted!70,minimum size=7.2mm,inner sep=0pt,font=\fontsize{8}{9}\selectfont] at (7,0) {8};
\draw[otMuted!80,line width=0.5pt] (6.72,0.28) -- (7.28,-0.28);
\node[draw=otRule,fill=white,text=otInk,minimum size=7.2mm,inner sep=0pt,font=\fontsize{8}{9}\selectfont] at (8,0) {9};
\node[draw=otRule,fill=otNavyTint!45,text=otMuted!70,minimum size=7.2mm,inner sep=0pt,font=\fontsize{8}{9}\selectfont] at (9,0) {10};
\draw[otMuted!80,line width=0.5pt] (8.72,0.28) -- (9.28,-0.28);
\node[draw=otRule,fill=white,text=otInk,minimum size=7.2mm,inner sep=0pt,font=\fontsize{8}{9}\selectfont] at (0,1) {11};
\node[draw=otRule,fill=otNavyTint!45,text=otMuted!70,minimum size=7.2mm,inner sep=0pt,font=\fontsize{8}{9}\selectfont] at (1,1) {12};
\draw[otMuted!80,line width=0.5pt] (0.72,1.28) -- (1.28,0.72);
\node[draw=otRule,fill=white,text=otInk,minimum size=7.2mm,inner sep=0pt,font=\fontsize{8}{9}\selectfont] at (2,1) {13};
\node[draw=otRule,fill=otNavyTint!45,text=otMuted!70,minimum size=7.2mm,inner sep=0pt,font=\fontsize{8}{9}\selectfont] at (3,1) {14};
\draw[otMuted!80,line width=0.5pt] (2.7199999999999998,1.28) -- (3.2800000000000002,0.72);
\node[draw=otRule,fill=white,text=otInk,minimum size=7.2mm,inner sep=0pt,font=\fontsize{8}{9}\selectfont] at (4,1) {15};
\node[draw=otRule,fill=otNavyTint!45,text=otMuted!70,minimum size=7.2mm,inner sep=0pt,font=\fontsize{8}{9}\selectfont] at (5,1) {16};
\draw[otMuted!80,line width=0.5pt] (4.72,1.28) -- (5.28,0.72);
\node[draw=otRule,fill=white,text=otInk,minimum size=7.2mm,inner sep=0pt,font=\fontsize{8}{9}\selectfont] at (6,1) {17};
\node[draw=otRule,fill=otNavyTint!45,text=otMuted!70,minimum size=7.2mm,inner sep=0pt,font=\fontsize{8}{9}\selectfont] at (7,1) {18};
\draw[otMuted!80,line width=0.5pt] (6.72,1.28) -- (7.28,0.72);
\node[draw=otRule,fill=white,text=otInk,minimum size=7.2mm,inner sep=0pt,font=\fontsize{8}{9}\selectfont] at (8,1) {19};
\node[draw=otRule,fill=otNavyTint!45,text=otMuted!70,minimum size=7.2mm,inner sep=0pt,font=\fontsize{8}{9}\selectfont] at (9,1) {20};
\draw[otMuted!80,line width=0.5pt] (8.72,1.28) -- (9.28,0.72);
\node[draw=otRule,fill=white,text=otInk,minimum size=7.2mm,inner sep=0pt,font=\fontsize{8}{9}\selectfont] at (0,2) {21};
\node[draw=otRule,fill=otNavyTint!45,text=otMuted!70,minimum size=7.2mm,inner sep=0pt,font=\fontsize{8}{9}\selectfont] at (1,2) {22};
\draw[otMuted!80,line width=0.5pt] (0.72,2.2800000000000002) -- (1.28,1.72);
\node[draw=otRule,fill=white,text=otInk,minimum size=7.2mm,inner sep=0pt,font=\fontsize{8}{9}\selectfont] at (2,2) {23};
\node[draw=otRule,fill=otNavyTint!45,text=otMuted!70,minimum size=7.2mm,inner sep=0pt,font=\fontsize{8}{9}\selectfont] at (3,2) {24};
\draw[otMuted!80,line width=0.5pt] (2.7199999999999998,2.2800000000000002) -- (3.2800000000000002,1.72);
\node[draw=otRule,fill=white,text=otInk,minimum size=7.2mm,inner sep=0pt,font=\fontsize{8}{9}\selectfont] at (4,2) {25};
\node[draw=otRule,fill=otNavyTint!45,text=otMuted!70,minimum size=7.2mm,inner sep=0pt,font=\fontsize{8}{9}\selectfont] at (5,2) {26};
\draw[otMuted!80,line width=0.5pt] (4.72,2.2800000000000002) -- (5.28,1.72);
\node[draw=otRule,fill=white,text=otInk,minimum size=7.2mm,inner sep=0pt,font=\fontsize{8}{9}\selectfont] at (6,2) {27};
\node[draw=otRule,fill=otNavyTint!45,text=otMuted!70,minimum size=7.2mm,inner sep=0pt,font=\fontsize{8}{9}\selectfont] at (7,2) {28};
\draw[otMuted!80,line width=0.5pt] (6.72,2.2800000000000002) -- (7.28,1.72);
\node[draw=otRule,fill=white,text=otInk,minimum size=7.2mm,inner sep=0pt,font=\fontsize{8}{9}\selectfont] at (8,2) {29};
\node[draw=otRule,fill=otNavyTint!45,text=otMuted!70,minimum size=7.2mm,inner sep=0pt,font=\fontsize{8}{9}\selectfont] at (9,2) {30};
\draw[otMuted!80,line width=0.5pt] (8.72,2.2800000000000002) -- (9.28,1.72);
\node[draw=otRule,fill=white,text=otInk,minimum size=7.2mm,inner sep=0pt,font=\fontsize{8}{9}\selectfont] at (0,3) {31};
\node[draw=otRule,fill=otNavyTint!45,text=otMuted!70,minimum size=7.2mm,inner sep=0pt,font=\fontsize{8}{9}\selectfont] at (1,3) {32};
\draw[otMuted!80,line width=0.5pt] (0.72,3.2800000000000002) -- (1.28,2.7199999999999998);
\node[draw=otRule,fill=white,text=otInk,minimum size=7.2mm,inner sep=0pt,font=\fontsize{8}{9}\selectfont] at (2,3) {33};
\node[draw=otRule,fill=otNavyTint!45,text=otMuted!70,minimum size=7.2mm,inner sep=0pt,font=\fontsize{8}{9}\selectfont] at (3,3) {34};
\draw[otMuted!80,line width=0.5pt] (2.7199999999999998,3.2800000000000002) -- (3.2800000000000002,2.7199999999999998);
\node[draw=otRule,fill=white,text=otInk,minimum size=7.2mm,inner sep=0pt,font=\fontsize{8}{9}\selectfont] at (4,3) {35};
\node[draw=otRule,fill=otNavyTint!45,text=otMuted!70,minimum size=7.2mm,inner sep=0pt,font=\fontsize{8}{9}\selectfont] at (5,3) {36};
\draw[otMuted!80,line width=0.5pt] (4.72,3.2800000000000002) -- (5.28,2.7199999999999998);
\node[draw=otRule,fill=white,text=otInk,minimum size=7.2mm,inner sep=0pt,font=\fontsize{8}{9}\selectfont] at (6,3) {37};
\node[draw=otRule,fill=otNavyTint!45,text=otMuted!70,minimum size=7.2mm,inner sep=0pt,font=\fontsize{8}{9}\selectfont] at (7,3) {38};
\draw[otMuted!80,line width=0.5pt] (6.72,3.2800000000000002) -- (7.28,2.7199999999999998);
\node[draw=otRule,fill=white,text=otInk,minimum size=7.2mm,inner sep=0pt,font=\fontsize{8}{9}\selectfont] at (8,3) {39};
\node[draw=otRule,fill=otNavyTint!45,text=otMuted!70,minimum size=7.2mm,inner sep=0pt,font=\fontsize{8}{9}\selectfont] at (9,3) {40};
\draw[otMuted!80,line width=0.5pt] (8.72,3.2800000000000002) -- (9.28,2.7199999999999998);
\node[draw=otRule,fill=white,text=otInk,minimum size=7.2mm,inner sep=0pt,font=\fontsize{8}{9}\selectfont] at (0,4) {41};
\node[draw=otRule,fill=otNavyTint!45,text=otMuted!70,minimum size=7.2mm,inner sep=0pt,font=\fontsize{8}{9}\selectfont] at (1,4) {42};
\draw[otMuted!80,line width=0.5pt] (0.72,4.28) -- (1.28,3.7199999999999998);
\node[draw=otRule,fill=white,text=otInk,minimum size=7.2mm,inner sep=0pt,font=\fontsize{8}{9}\selectfont] at (2,4) {43};
\node[draw=otRule,fill=otNavyTint!45,text=otMuted!70,minimum size=7.2mm,inner sep=0pt,font=\fontsize{8}{9}\selectfont] at (3,4) {44};
\draw[otMuted!80,line width=0.5pt] (2.7199999999999998,4.28) -- (3.2800000000000002,3.7199999999999998);
\node[draw=otRule,fill=white,text=otInk,minimum size=7.2mm,inner sep=0pt,font=\fontsize{8}{9}\selectfont] at (4,4) {45};
\node[draw=otRule,fill=otNavyTint!45,text=otMuted!70,minimum size=7.2mm,inner sep=0pt,font=\fontsize{8}{9}\selectfont] at (5,4) {46};
\draw[otMuted!80,line width=0.5pt] (4.72,4.28) -- (5.28,3.7199999999999998);
\node[draw=otRule,fill=white,text=otInk,minimum size=7.2mm,inner sep=0pt,font=\fontsize{8}{9}\selectfont] at (6,4) {47};
\node[draw=otRule,fill=otNavyTint!45,text=otMuted!70,minimum size=7.2mm,inner sep=0pt,font=\fontsize{8}{9}\selectfont] at (7,4) {48};
\draw[otMuted!80,line width=0.5pt] (6.72,4.28) -- (7.28,3.7199999999999998);
\node[draw=otRule,fill=white,text=otInk,minimum size=7.2mm,inner sep=0pt,font=\fontsize{8}{9}\selectfont] at (8,4) {49};
\node[draw=otRule,fill=otNavyTint!45,text=otMuted!70,minimum size=7.2mm,inner sep=0pt,font=\fontsize{8}{9}\selectfont] at (9,4) {50};
\draw[otMuted!80,line width=0.5pt] (8.72,4.28) -- (9.28,3.7199999999999998);
\end{tikzpicture}}\\[1mm]
\legende{Après le premier tour ($i=2$) : les multiples de $2$ sont barrés.}
\end{codecote}
\end{exercice}

\begin{exercice}[3]{Simuler la somme de deux dés}
On lance deux dés équilibrés et on note $S$ la somme obtenue. Le programme suivant simule $n=10\,000$ lancers.
\begin{python}
from random import randint

effectifs = [0] * 13
n = 10000
for k in range(n):
    s = randint(1, 6) + randint(1, 6)
    effectifs[s] = effectifs[s] + 1
frequences = [effectifs[s] / n for s in range(2, 13)]
\end{python}
\code{randint(1, 6)} renvoie un entier au hasard entre $1$ et $6$ (inclus). Une exécution a donné les effectifs suivants.\par
\nopagebreak\smallskip
{\renewcommand{\arraystretch}{1.3}\arrayrulecolor{otRule}%
\begin{tabular*}{\linewidth}{@{\extracolsep{\fill}}lccccccccccc@{}}
\toprule
\otTete{Somme $s$} & $2$ & $3$ & $4$ & $5$ & $6$ & $7$ & $8$ & $9$ & $10$ & $11$ & $12$ \\
\midrule
\otTete{Effectif} & $274$ & $591$ & $841$ & $1\,214$ & $1\,354$ & $1\,642$ & $1\,354$ & $1\,090$ & $801$ & $533$ & $306$ \\
\bottomrule
\end{tabular*}}
\begin{questions}
  \item Pourquoi la liste \code{effectifs} a-t-elle $13$ éléments ? À la fin du programme, que valent \code{effectifs[0]} et
    \code{effectifs[1]} ?
  \item Justifier que, pour $s$ entier de $2$ à $12$, $P(S=s)=\frac{6-|s-7|}{36}$.
  \item Calculer les fréquences obtenues et les comparer aux probabilités. Quelle est la somme la plus fréquente ?
  \item À l'aide de l'inégalité de Bienaymé-Tchebychev, majorer la probabilité que la fréquence $F$ de la somme $7$ sur $10\,000$
    lancers s'écarte de $\frac16$ d'au moins $0{,}01$. Le résultat de la simulation est-il surprenant ?
\end{questions}
\end{exercice}
